% Vietnamese translation for OI Public Library
% Source-File: IOI中国国家候选队论文/2006/朱晨光/朱晨光.ppt
\documentclass[11pt]{article}
\usepackage{oipl-vn}

\newcommand{\Oh}{\mathcal{O}}

\title{Ứng dụng của cấu trúc dữ liệu cơ bản\\{\large Ghi chú theo slide}}
\author{Zhu Chenguang}
\date{2006}

\begin{document}
\maketitle

\origtitle{Applications of basic data structures}
\sourcepath{IOI China candidate team papers / 2006 / Zhu Chenguang / Zhu Chenguang.ppt}

\section{Phạm vi}

Bộ trình chiếu đi kèm bài viết về danh sách tuyến tính, ngăn xếp và hàng đợi
trong thi tin học. Ghi chú này rút lại các ý chính từ bản bài viết đã dịch.

\section{Cấu trúc cơ bản vẫn quan trọng}

Các cấu trúc nâng cao không phải lúc nào cũng là lựa chọn tốt nhất trong phòng
thi. Mảng, danh sách liên kết, ngăn xếp và hàng đợi có chi phí cài đặt thấp,
dễ kiểm tra, và thường đủ để khai thác đúng tính chất của đề.

\section{Ngăn xếp}

Ngăn xếp phù hợp với quan hệ vào sau ra trước: khử ngoặc, duyệt sâu, mô phỏng
đệ quy, duy trì dãy đơn điệu, hoặc xử lý các phần tử đang chờ một biên tiếp
theo. Điểm cần chú ý là phần tử bị lấy ra khỏi ngăn xếp thường đã được chứng
minh không còn cần cho tương lai.

\section{Danh sách tuyến tính}

Danh sách tuyến tính có thể biểu diễn bằng mảng hoặc liên kết. Mảng cho truy
cập theo chỉ số trong \(\Oh(1)\), còn danh sách liên kết thuận tiện khi cần
xóa, chèn nhiều phần tử đã biết vị trí. Trong nhiều bài, một danh sách đôi mô
phỏng đúng quá trình biến đổi mà không cần cây cân bằng.

\section{Hàng đợi}

Hàng đợi tự nhiên cho tìm kiếm theo chiều rộng và các quá trình lan truyền.
Khi thêm điều kiện đơn điệu, nó trở thành hàng đợi đơn điệu, một công cụ mạnh
để tối ưu quy hoạch động hoặc truy vấn cửa sổ trượt.

\section{Ghi nhớ}

Trước khi dùng cấu trúc phức tạp, nên hỏi cấu trúc dữ liệu tối thiểu cần hỗ
trợ những thao tác nào. Nếu các thao tác chỉ là thêm cuối, lấy đầu, xóa phần
tử kề, hoặc quay lui, một cấu trúc cơ bản thường đem lại lời giải gọn và ít
rủi ro hơn.

\end{document}
